% Ad Atticum 2.1 --- headnote
% ad-atticum-02-01, 2 or 3 June 60 BC, written at Rome (or just outside)

\textit{Cicero to Atticus, written at Rome on or just after
the Kalends of June 60 BC. The first letter of book 2 of
Atticus, opening on the dating-line that begins the political
year of 60: ``On the Kalends of June, as I was setting out
for Antium and gladly leaving Marcus Metellus's gladiators
behind, your boy met me.''}

\textit{Three large pieces. First, the literature
(\textsection 1--3): Atticus has sent a Greek memoir of
Cicero's consulship; Cicero has just sent his own Greek
memoir to Atticus through Cossinius. Each had written without
knowing the other's would arrive; Cicero finds Atticus's
``a little rough and uncombed,'' a kind of female-without-
perfume Atticism, while his own ``has used up the whole
perfumery-box of Isocrates and all his pupils' little
caskets, and not a few of Aristotle's pigments too.''
Posidonius at Rhodes, sent the memoir to expand into a fuller
work, has been ``not stirred to writing but plainly
deterred.'' Cicero offers the famous list of his ten
``consular speeches'' --- the canonical corpus he has been
publishing as a body, after the model of Demosthenes's
Philippics: \textit{De Lege Agraria} I (in the senate, Kal.
Jan.), II (to the people, on the agrarian law), III (on
Otho), \textit{Pro Rabirio}, the lost speech \textit{de
proscriptorum filiis}, the senate-speech in which he gave
up his province, the four \textit{Catilinarians}, plus the
two short fragments of agrarian-law oratory; the corpus
later becomes the standard ``consular orations'' edition.}

\textit{Second, Clodius (\textsection 4--5). Pulchellus
``the pretty boy'' is now openly canvassing for the tribunate
of the plebs as a transferred plebeian. Cicero, in the senate,
has destroyed the boasts about how fast Clodius travels (``from
Sicily to Rome in seven days; three hours before that, from
Rome to Interamna; you entered by night --- the same as
before; no one came out to meet you --- not even then, when
they most ought to have come''). At a candidate's escort,
Clodius's joke about his sister Clodia ``giving him only one
foot of room'' provokes Cicero's famous reply: ``Don't
complain of one foot of your sister's; you are permitted to
lift the other one too.'' The \textit{not a consular saying}
the editor admits and defends.}

\textit{Third, the Pompey-Caesar policy (\textsection 6--8).
Cicero defends his now-close partnership with Pompey, and adds
that he means to draw Caesar (whose ``winds are now greatly
favourable'') in the same direction. The opposing voice in his
own party is Cato, who ``with the best mind and the highest
faith sometimes harms the commonwealth; he gives his opinion
as if he were in the Republic of Plato, not in the dregs of
Romulus.'' The image of the rich \textit{piscinarii} content
with bearded mullets that come to the hand returns. Closing
business: the avenger-of-debt buried in debt, the books of
Servius Claudius secured through Paetus, a quiet courtesy to
Octavius about Atticus's provincial affairs.}
